\documentclass[12pt]{article}

\usepackage[a4paper,margin=1in]{geometry}
\usepackage{amsmath,amssymb}
\usepackage{newtxtext,newtxmath}
\usepackage{fancyhdr}
\usepackage{listings}
\usepackage{graphicx}
\usepackage{float}
\usepackage{hyperref}
\usepackage{booktabs}
\usepackage{xcolor}

\setlength{\parindent}{0pt}
\setlength{\parskip}{5pt}

% Keep entire document at 12 pt
\AtBeginDocument{
    \fontsize{12pt}{14pt}\selectfont
}

% ------------------------------------------------
% Header Configuration
% ------------------------------------------------
\pagestyle{fancy}
\fancyhf{}
\fancyhead[L]{Name: \underline{Shaik Jnaved}}
\fancyhead[R]{Enrollment No.: \underline{P26DS001}}
\renewcommand{\headrulewidth}{0.4pt}
\setlength{\headheight}{15pt}

% ------------------------------------------------
% Code Formatting Configuration
% ------------------------------------------------
\lstset{
    basicstyle=\ttfamily\fontsize{8.5pt}{10.5pt}\selectfont,
    breaklines=true,
    frame=single,
    columns=fullflexible,
    keepspaces=true,
    showstringspaces=false,
    keywordstyle=\color{blue}\bfseries,
    commentstyle=\color{gray}\itshape,
    stringstyle=\color{teal}
}

\begin{document}

% ==============================================================================
% TITLE BANNER
% ==============================================================================
\begin{center}
\textbf{\Large Computer Science and Engineering Department, SVNIT, Surat}\\[4pt]
\textbf{\large M.Tech. I -- Semester I}\\[3pt]
\textbf{\large Design and Analysis of Algorithm (CSDS103)}\\[4pt]
\textbf{\Large Lab Assignment 7: String Matching Algorithms \& Empirical Analysis}
\end{center}

\vspace{0.2cm}
\hrule
\vspace{0.4cm}

\textbf{Abstract}:\\
String matching is a foundational algorithmic problem in computer science with applications ranging from text editors, web search engines, and network packet inspection to computational genomics and bioinformatics. Given a text string $T$ of length $n$ and a pattern string $P$ of length $m$ (with $m \le n$), the objective is to determine all valid shift positions $s \in \{0, 1, \dots, n-m\}$ such that $T[s \dots s+m-1] = P[0 \dots m-1]$. 

This laboratory report presents an end-to-end mathematical, algorithmic, and empirical investigation of three canonical string matching algorithms:
\begin{enumerate}
    \item \textbf{Question 1: The Naive (Brute-Force) String Matching Approach}
    \item \textbf{Question 2: The Knuth-Morris-Pratt (KMP) Linear-Time Algorithm}
    \item \textbf{Question 3: The Rabin-Karp Rolling Hash Matching Algorithm}
    \item \textbf{Question 4: Comprehensive Empirical Benchmark \& Performance Analysis}
\end{enumerate}
We evaluate the operational behavior, theoretical complexities, empirical execution times across variable input sizes and pattern lengths, and the actual number of character comparisons incurred by each method.

\vspace{0.4cm}

% ==============================================================================
% QUESTION 1: NAIVE STRING MATCHING
% ==============================================================================
\section*{Question 1: Naive (Brute-Force) String Matching Approach}

\subsection*{1.1 Problem Statement and Algorithmic Principle}
The Naive string matching algorithm is the simplest approach to finding occurrences of pattern $P[0 \dots m-1]$ in text $T[0 \dots n-1]$. It operates by sliding a window of length $m$ over the text $T$, sequentially testing every possible starting position (shift) $s \in \{0, 1, \dots, n-m\}$.

For each shift $s$:
\begin{itemize}
    \item It aligns $P[0]$ with $T[s]$ and compares characters $T[s+j]$ with $P[j]$ for $j = 0, 1, \dots, m-1$.
    \item If a character mismatch is detected at index $j$ ($T[s+j] \ne P[j]$), the inner loop terminates immediately, the shift is declared invalid, and the search proceeds to shift $s+1$.
    \item If all $m$ characters match successfully, the valid shift index $s$ is appended to the results list.
\end{itemize}

A fundamental limitation of the Naive algorithm is that it discards all historical character match information upon encountering a mismatch, often re-examining characters of $T$ that were already inspected during previous shifts.

\subsection*{1.2 Python Implementation (\texttt{naive\_string\_match.py})}
The source code implemented in \texttt{src/lab\_7/naive\_string\_match.py} is presented below:

\begin{lstlisting}[language=Python]
def search(text: str, pattern: str):
    n = len(text)
    m = len(pattern)
    res = []

    for i in range(n-m+1):
        k = i
        flag = True
        for j in range(m):
            if text[k] == pattern[j]:
                j+=1
                k+=1
            else:
                flag = False
                break
        
        if flag:
            res.append(i)
    
    return res


def main(text: str, pattern: str):
    search(text=text, pattern=pattern)


if __name__ == "__main__":
    text = "aabaacaadaabaaba"
    pattern = "aaba"

    res = search(text=text, pattern=pattern)
    print(res)
\end{lstlisting}

\subsection*{1.3 Theoretical Complexity Analysis}
\begin{itemize}
    \item \textbf{Worst-Case Time Complexity $\mathcal{O}((n - m + 1) \cdot m) = \mathcal{O}(n \cdot m)$}:\\
    Occurs when the pattern and text have extensive repetitions with mismatches occurring only at the very end of the pattern. For instance, if $T = \text{"aaaa...a"}$ ($n$ times 'a') and $P = \text{"aaaab"}$ ($m-1$ times 'a' followed by 'b'), every shift checks $m$ characters before failing, leading to $(n - m + 1) \cdot m$ total character comparisons.
    
    \item \textbf{Best-Case Time Complexity $\mathcal{O}(n - m + 1) = \mathcal{O}(n)$}:\\
    Occurs when the very first character of the pattern never matches the text at any shift ($T[s] \ne P[0]$). Only 1 comparison is performed per shift, yielding exactly $n - m + 1$ comparisons.
    
    \item \textbf{Average-Case Time Complexity $\mathcal{O}(n)$}:\\
    For random text over an alphabet of size $|\Sigma| = d$, the probability that two random characters match is $1/d$. The expected number of comparisons at each shift is a geometric series:
    \[
    \mathbb{E}[\text{comps per shift}] = \sum_{k=1}^{m} \left(\frac{1}{d}\right)^{k-1} \le \frac{d}{d-1} = \mathcal{O}(1)
    \]
    Thus, across all $n-m+1$ shifts, the expected number of comparisons is linear in $n$.
    
    \item \textbf{Auxiliary Space Complexity $\mathcal{O}(1)$}:\\
    The algorithm requires only a few pointer variables ($i, j, k, n, m$) and no auxiliary data structures beyond storing the result indices.
\end{itemize}

\subsection*{1.4 Verification \& Sample Output}
Running the sample execution on test text $T = \text{"aabaacaadaabaaba"}$ and pattern $P = \text{"aaba"}$:
\begin{itemize}
    \item Shift $s = 0$: $T[0\dots 3] = \text{"aaba"}$ matches $P$ $\implies$ record index $0$.
    \item Shift $s = 9$: $T[9\dots 12] = \text{"aaba"}$ matches $P$ $\implies$ record index $9$.
    \item Shift $s = 12$: $T[12\dots 15] = \text{"aaba"}$ matches $P$ $\implies$ record index $12$.
\end{itemize}
\textbf{Output}: \texttt{[0, 9, 12]}. Correctness verified.

\vspace{0.4cm}

% ==============================================================================
% QUESTION 2: KMP ALGORITHM
% ==============================================================================
\section*{Question 2: Knuth-Morris-Pratt (KMP) Algorithm}

\subsection*{2.1 Problem Statement and Algorithmic Principle}
The Knuth-Morris-Pratt (KMP) algorithm achieves deterministic worst-case linear time $\Theta(n + m)$ by ensuring that the text pointer $i$ never moves backward (backtracks). It exploits the internal structure of the pattern itself through a precomputed \textbf{Longest Proper Prefix which is also a Suffix (LPS)} array, commonly referred to as the $\pi$ table.

\textbf{1. LPS Array Definition}:
For a pattern $P$ of length $m$, $\text{lps}[j]$ stores the length of the longest proper prefix of $P[0 \dots j]$ that is also a suffix of $P[0 \dots j]$. 
\begin{itemize}
    \item A \textit{proper prefix} of a string $S$ is a prefix that is not equal to $S$ itself.
    \item $\text{lps}[0] = 0$ by definition.
\end{itemize}

\textbf{2. Preprocessing Phase ($\Theta(m)$)}:
The LPS array is constructed iteratively using two pointers (\texttt{i} and \texttt{pattern\_length}). If $P[i] == P[\text{pattern\_length}]$, we increment \texttt{pattern\_length} and store it in $\text{lps}[i]$. If they mismatch and $\text{pattern\_length} > 0$, \texttt{pattern\_length} falls back to $\text{lps}[\text{pattern\_length}-1]$ without resetting to zero.

\textbf{3. Search Phase ($\Theta(n)$)}:
During matching, we maintain text index $i$ and pattern index $j$:
\begin{itemize}
    \item When $T[i] == P[j]$, both $i$ and $j$ advance.
    \item If a complete match occurs ($j == m$), index $i - j$ is recorded as a match, and $j$ is updated to $\text{lps}[j-1]$ to find subsequent overlapping occurrences.
    \item If a mismatch occurs after matching $j$ characters ($T[i] \ne P[j]$):
    \begin{itemize}
        \item If $j > 0$, we slide the pattern to align its longest matched prefix with the matched suffix of the text by setting $j = \text{lps}[j-1]$. Pointer $i$ remains unchanged!
        \item If $j == 0$, no prefix could match, so we simply increment $i$.
    \end{itemize}
\end{itemize}

\subsection*{2.2 Python Implementation (\texttt{kmp.py})}
The source code implemented in \texttt{src/lab\_7/kmp.py} is presented below:

\begin{lstlisting}[language=Python]
def construct_lps(pattern: str, lps: list):
    m = len(pattern)
    pattern_length = 0

    lps[0] = 0
    i=1
    while(i<m):
        if pattern[i] == pattern[pattern_length]:
            pattern_length += 1
            lps[i] = pattern_length
            i+=1
        else:
            if pattern_length != 0:
                pattern_length = lps[pattern_length-1]
            else:
                lps[i] = 0
                i+=1


def search(text: str, pattern: str):
    n = len(text)
    m = len(pattern)

    # lps -> longest proper prefix which is also a suffix
    lps = [0]*m
    res = []

    construct_lps(pattern=pattern, lps=lps)
    i, j = 0, 0

    while(i<n):
        if text[i] == pattern[j]:
            i+=1
            j+=1

            if j==m:
                res.append(i-j)
                j = lps[j-1]
        else:
            if j!=0:
                j = lps[j-1]
            else:
                i+=1
    return res


def main(text: str, pattern: str):
    search(text=text, pattern=pattern)


if __name__ == "__main__":
    text = "aabaacaadaabaaba"
    pattern = "aaba"
    res = search(text=text, pattern=pattern)
    print(res)
\end{lstlisting}

\subsection*{2.3 Theoretical Complexity Analysis}
\begin{itemize}
    \item \textbf{Preprocessing Time Complexity $\Theta(m)$}:\\
    In \texttt{construct\_lps}, variable $i$ increments at each step when characters match or when $\text{pattern\_length} == 0$. The potential function $\Phi = \text{pattern\_length} \ge 0$ increases by 1 only when $i$ increases and strictly decreases during fallback. Thus, the total number of operations is bounded by $2m = \Theta(m)$.

    \item \textbf{Search Time Complexity $\Theta(n)$}:\\
    Similarly, in \texttt{search}, pointer $i$ strictly increases from $0$ to $n$. The pattern pointer $j$ increases at most $n$ times (once per match). Every fallback $j = \text{lps}[j-1]$ strictly decreases $j$ without decreasing $i$. Since $j \ge 0$, $j$ can decrease at most as many times as it increased. Hence, the while loop executes at most $2n$ times.

    \item \textbf{Total Time Complexity $\Theta(n + m)$}:\\
    The algorithm is strictly linear in all cases (best, average, and worst case), providing robust guarantees against pathological inputs.

    \item \textbf{Auxiliary Space Complexity $\Theta(m)$}:\\
    An auxiliary integer array of size $m$ is required to store the LPS table.
\end{itemize}

\subsection*{2.4 Verification \& Sample Output}
For pattern $P = \text{"aaba"}$:
\begin{itemize}
    \item $P[0] = \text{'a'} \implies \text{lps}[0] = 0$.
    \item $P[1] = \text{'a'} == P[0] \implies \text{lps}[1] = 1$.
    \item $P[2] = \text{'b'} \ne P[1]$, fallback to $P[\text{lps}[0]] = P[0] = \text{'a'} \ne \text{'b'} \implies \text{lps}[2] = 0$.
    \item $P[3] = \text{'a'} == P[0] \implies \text{lps}[3] = 1$.
    \item LPS array: \texttt{[0, 1, 0, 1]}.
\end{itemize}
Searching on $T = \text{"aabaacaadaabaaba"}$ yields match indices: \texttt{[0, 9, 12]}. Correctness verified.

\vspace{0.4cm}

% ==============================================================================
% QUESTION 3: RABIN-KARP ALGORITHM
% ==============================================================================
\section*{Question 3: Rabin-Karp String Matching Algorithm}

\subsection*{3.1 Problem Statement and Algorithmic Principle}
The Rabin-Karp algorithm introduces a hashing-based approach to string matching. Rather than performing character-by-character checks at every shift position, it associates numerical hash values with the pattern and with each window of text.

\textbf{1. Rolling Hash Principle}:\\
Treating strings over an alphabet as base-$d$ numbers (in our implementation, $d = 10$) modulo a selected prime $Q = 101$:
\begin{itemize}
    \item The pattern hash is:
    \[
    p = \left( \sum_{i=0}^{m-1} P[i] \cdot d^{m-1-i} \right) \bmod Q
    \]
    \item The hash of the initial text window $T[0 \dots m-1]$ is computed identically as $t_0$.
\end{itemize}

\textbf{2. $\mathcal{O}(1)$ Sliding Window Hash Update (Horner's Rule)}:\\
When sliding the window from shift $s$ to $s+1$, the old high-order character $T[s]$ is subtracted, the hash is multiplied by base $d$, and the new incoming character $T[s+m]$ is added:
\[
t_{s+1} = \left( d \cdot (t_s - T[s] \cdot h) + T[s+m] \right) \bmod Q
\]
where $h = d^{m-1} \bmod Q$ is precalculated in $\mathcal{O}(\log m)$ using modular exponentiation. If the subtraction yields a negative result, $Q$ is added to restore modular positivity. This rolling update executes in $\mathcal{O}(1)$ arithmetic operations.

\textbf{3. Spurious Hit Handling (Verification)}:\\
If $t_s == p$, two possibilities exist:
\begin{enumerate}
    \item \textbf{True Match}: $T[s \dots s+m-1] == P$.
    \item \textbf{Spurious Hit (Hash Collision)}: $t_s \equiv p \pmod Q$, but $T[s \dots s+m-1] \ne P$.
\end{enumerate}
To guarantee accuracy, whenever $t_s == p$, a character-by-character verification $T[s \dots s+m-1] == P$ is performed.

\subsection*{3.2 Python Implementation (\texttt{rabin\_karp.py})}
The source code implemented in \texttt{src/lab\_7/rabin\_karp.py} is presented below:

\begin{lstlisting}[language=Python]
Q = 101


def search(text: str, pattern: str):
    n = len(text)
    m = len(pattern)

    res = []
    h = pow(10, m-1, Q)

    p_hash, t_hash = 0, 0

    # initial hash
    for i in range(m):
        p_hash = (10 * p_hash + ord(pattern[i])) % Q
        t_hash = (10 * t_hash + ord(text[i])) % Q

    for i in range(n-m+1):
        if p_hash == t_hash:
            if text[i: i+m] == pattern:
                res.append(i)

        if i < n-m:
            t_hash = (10 * (t_hash - ord(text[i]) * h) + ord(text[i+m])) % Q
            if t_hash < 0:
                t_hash += Q

    return res


def main(text: str, pattern: str):
    search(text=text, pattern=pattern)


if __name__ == "__main__":
    text = "aabaacaadaabaaba"
    pattern = "aaba"
    res = search(text=text, pattern=pattern)
    print(res)
\end{lstlisting}

\subsection*{3.3 Theoretical Complexity Analysis}
\begin{itemize}
    \item \textbf{Preprocessing Time Complexity $\Theta(m)$}:\\
    Computing $h = d^{m-1} \bmod Q$ and evaluating the initial hashes $p$ and $t_0$ takes $\Theta(m)$ arithmetic steps.

    \item \textbf{Average / Best-Case Time Complexity $\mathcal{O}(n + m)$}:\\
    If prime $Q$ is chosen sufficiently large relative to $m$, the probability of spurious hits is approximately $1/Q$. If the expected number of collisions is $\mathcal{O}(1)$, verification is triggered only $\mathcal{O}(1)$ times. Rolling hash updates take $\mathcal{O}(1)$ per shift for $n-m$ shifts, yielding total average time $\mathcal{O}(n + m)$.

    \item \textbf{Worst-Case Time Complexity $\mathcal{O}((n - m + 1) \cdot m) = \mathcal{O}(n \cdot m)$}:\\
    Occurs when every shift produces a hash collision ($t_s == p$) that fails during verification, or when the pattern occurs at every position. In this degenerate case, all $m$ characters are verified at every shift.

    \item \textbf{Auxiliary Space Complexity $\mathcal{O}(1)$}:\\
    Rabin-Karp requires only a constant number of integer variables ($p\_hash, t\_hash, h, Q$) for maintaining hash states.
\end{itemize}

\subsection*{3.4 Verification \& Sample Output}
Running the algorithm on $T = \text{"aabaacaadaabaaba"}$ and $P = \text{"aaba"}$ with prime $Q = 101$:
\begin{itemize}
    \item Hash matches occur at shifts $0, 9, 12$ (verified as true matches).
    \item Output: \texttt{[0, 9, 12]}. Correctness verified.
\end{itemize}

\vspace{0.4cm}

% ==============================================================================
% QUESTION 4: EMPIRICAL PERFORMANCE ANALYSIS
% ==============================================================================
\section*{Question 4: Experimental Performance Analysis \& Empirical Comparisons}

\subsection*{4.1 Experimental Methodology and Setup}
To provide a thorough empirical assessment, we developed a comprehensive benchmarking framework in \texttt{src/lab\_7/comparisions.py}. The testing environment and parameters are structured as follows:
\begin{itemize}
    \item \textbf{Alphabet}: Controlled alphabet $\Sigma = \{\text{'a'}, \text{'b'}, \text{'c'}, \text{'d'}\}$ ($|\Sigma| = 4$).
    \item \textbf{Text Sizes}: $n \in \{500, 1000, 10000\}$ characters, randomly sampled with reproducible seeds.
    \item \textbf{Pattern Lengths}: $m \in \{5, 10, 20\}$ characters, sliced directly from text to ensure valid matches exist.
    \item \textbf{Timing Metric}: Measured using high-resolution \texttt{time.perf\_counter()}, with results averaged over multiple repetitions to eliminate OS thread-scheduling anomalies.
    \item \textbf{Comparison Counter}: Dedicated character comparison functions tracking actual equality checks performed during execution.
\end{itemize}

\subsection*{4.2 Source Code (\texttt{comparisions.py})}
The complete driver script implemented in \texttt{src/lab\_7/comparisions.py} is shown below:

\begin{lstlisting}[language=Python]
import time
import random
from pathlib import Path
import matplotlib.pyplot as plt
from naive_string_match import main as naive
from kmp import main as kmp
from rabin_karp import main as rk

CURRENT_DIR = Path(__file__).resolve().parent
IMAGES_DIR = CURRENT_DIR.parent.parent / "images"
IMAGES_DIR.mkdir(parents=True, exist_ok=True)

def generate_random_text(length: int | str, alphabet: str = "abcd", seed: int | None = None) -> str:
    rng = random.Random(seed) if seed is not None else random
    return "".join(rng.choices(alphabet, k=int(length)))

def generate_random_pattern(length: int | str, alphabet: str = "abcd", text: str | None = None, seed: int | None = None) -> str:
    rng = random.Random(seed) if seed is not None else random
    length = int(length)
    if text is not None and len(text) >= length:
        start_idx = rng.randint(0, len(text) - length)
        return text[start_idx : start_idx + length]
    return "".join(rng.choices(alphabet, k=length))

def count_naive_comparisons(text: str, pattern: str) -> int:
    n, m = len(text), len(pattern)
    comps = 0
    for i in range(n - m + 1):
        k = i
        for j in range(m):
            comps += 1
            if text[k] == pattern[j]:
                k += 1
            else:
                break
    return comps

def count_kmp_comparisons(text: str, pattern: str) -> int:
    m, n = len(pattern), len(text)
    lps = [0] * m
    pattern_length, i, comps = 0, 1, 0
    while i < m:
        comps += 1
        if pattern[i] == pattern[pattern_length]:
            pattern_length += 1
            lps[i] = pattern_length
            i += 1
        else:
            if pattern_length != 0:
                pattern_length = lps[pattern_length - 1]
            else:
                lps[i] = 0
                i += 1
    i, j = 0, 0
    while i < n:
        comps += 1
        if text[i] == pattern[j]:
            i += 1
            j += 1
            if j == m:
                j = lps[j - 1]
        else:
            if j != 0:
                j = lps[j - 1]
            else:
                i += 1
    return comps

def count_rk_comparisons(text: str, pattern: str, Q: int = 101) -> int:
    n, m = len(text), len(pattern)
    h = pow(10, m - 1, Q)
    p_hash, t_hash, comps = 0, 0, 0
    for i in range(m):
        p_hash = (10 * p_hash + ord(pattern[i])) % Q
        t_hash = (10 * t_hash + ord(text[i])) % Q
    for i in range(n - m + 1):
        if p_hash == t_hash:
            for j in range(m):
                comps += 1
                if text[i + j] != pattern[j]:
                    break
        if i < n - m:
            t_hash = (10 * (t_hash - ord(text[i]) * h) + ord(text[i + m])) % Q
            if t_hash < 0:
                t_hash += Q
    return comps

input_sizes = [500, 1000, 10000]
pattern_sizes = [5, 10, 20]
n = len(input_sizes)
texts = [generate_random_text(length=i, seed=42 + idx) for idx, i in enumerate(input_sizes)]
patterns = [generate_random_pattern(length=pattern_sizes[i], text=texts[i], seed=100 + i) for i in range(n)]

def input_size_vs_execution_time(show_plot: bool = True):
    naive_map, kmp_map, rk_map = {}, {}, {}
    iterations = 5
    for i in range(n):
        t = time.perf_counter()
        for _ in range(iterations): naive(text=texts[i], pattern=patterns[i])
        naive_map[input_sizes[i]] = ((time.perf_counter() - t) / iterations) * 1000

        t = time.perf_counter()
        for _ in range(iterations): kmp(text=texts[i], pattern=patterns[i])
        kmp_map[input_sizes[i]] = ((time.perf_counter() - t) / iterations) * 1000

        t = time.perf_counter()
        for _ in range(iterations): rk(text=texts[i], pattern=patterns[i])
        rk_map[input_sizes[i]] = ((time.perf_counter() - t) / iterations) * 1000

    plt.figure(figsize=(9, 5.5))
    plt.plot(input_sizes, [naive_map[s] for s in input_sizes], label="Naive String Match", marker="o", color="#d62728", linewidth=2)
    plt.plot(input_sizes, [kmp_map[s] for s in input_sizes], label="KMP", marker="s", color="#1f77b4", linewidth=2)
    plt.plot(input_sizes, [rk_map[s] for s in input_sizes], label="Rabin-Karp", marker="^", color="#2ca02c", linewidth=2)
    plt.title("Input Size vs. Execution Time", fontsize=13, fontweight="bold")
    plt.xlabel("Text Size (n)", fontsize=11)
    plt.ylabel("Execution Time (ms)", fontsize=11)
    plt.grid(True, linestyle="--", alpha=0.6)
    plt.legend(fontsize=10)
    plt.tight_layout()
    plt.savefig(IMAGES_DIR / "input_size_vs_execution_time.png", dpi=300)
    if show_plot: plt.show()

def pattern_length_vs_execution_time(show_plot: bool = True):
    naive_map, kmp_map, rk_map = {}, {}, {}
    iterations = 5
    fixed_text = texts[-1]
    test_patterns = [generate_random_pattern(length=ps, text=fixed_text, seed=200 + idx) for idx, ps in enumerate(pattern_sizes)]
    for i in range(len(pattern_sizes)):
        ps, p = pattern_sizes[i], test_patterns[i]
        t = time.perf_counter()
        for _ in range(iterations): naive(text=fixed_text, pattern=p)
        naive_map[ps] = ((time.perf_counter() - t) / iterations) * 1000
        t = time.perf_counter()
        for _ in range(iterations): kmp(text=fixed_text, pattern=p)
        kmp_map[ps] = ((time.perf_counter() - t) / iterations) * 1000
        t = time.perf_counter()
        for _ in range(iterations): rk(text=fixed_text, pattern=p)
        rk_map[ps] = ((time.perf_counter() - t) / iterations) * 1000

    plt.figure(figsize=(9, 5.5))
    plt.plot(pattern_sizes, [naive_map[p] for p in pattern_sizes], label="Naive String Match", marker="o", color="#d62728", linewidth=2)
    plt.plot(pattern_sizes, [kmp_map[p] for p in pattern_sizes], label="KMP", marker="s", color="#1f77b4", linewidth=2)
    plt.plot(pattern_sizes, [rk_map[p] for p in pattern_sizes], label="Rabin-Karp", marker="^", color="#2ca02c", linewidth=2)
    plt.title(f"Pattern Length vs. Execution Time (Text Size = {len(fixed_text)})", fontsize=13, fontweight="bold")
    plt.xlabel("Pattern Length (m)", fontsize=11)
    plt.ylabel("Execution Time (ms)", fontsize=11)
    plt.grid(True, linestyle="--", alpha=0.6)
    plt.legend(fontsize=10)
    plt.tight_layout()
    plt.savefig(IMAGES_DIR / "pattern_length_vs_execution_time.png", dpi=300)
    if show_plot: plt.show()

def input_size_vs_comparisions(show_plot: bool = True):
    naive_map, kmp_map, rk_map = {}, {}, {}
    for i in range(n):
        sz = input_sizes[i]
        naive_map[sz] = count_naive_comparisons(texts[i], patterns[i])
        kmp_map[sz] = count_kmp_comparisons(texts[i], patterns[i])
        rk_map[sz] = count_rk_comparisons(texts[i], patterns[i])

    plt.figure(figsize=(9, 5.5))
    plt.plot(input_sizes, [naive_map[s] for s in input_sizes], label="Naive String Match", marker="o", color="#d62728", linewidth=2)
    plt.plot(input_sizes, [kmp_map[s] for s in input_sizes], label="KMP", marker="s", color="#1f77b4", linewidth=2)
    plt.plot(input_sizes, [rk_map[s] for s in input_sizes], label="Rabin-Karp (Verification Comps)", marker="^", color="#2ca02c", linewidth=2)
    plt.title("Input Size vs. Number of Comparisons", fontsize=13, fontweight="bold")
    plt.xlabel("Text Size (n)", fontsize=11)
    plt.ylabel("Number of Character Comparisons", fontsize=11)
    plt.grid(True, linestyle="--", alpha=0.6)
    plt.legend(fontsize=10)
    plt.tight_layout()
    plt.savefig(IMAGES_DIR / "input_size_vs_comparisons.png", dpi=300)
    if show_plot: plt.show()

if __name__ == "__main__":
    input_size_vs_execution_time()
    pattern_length_vs_execution_time()
    input_size_vs_comparisions()
\end{lstlisting}

\newpage

\subsection*{4.3 Analysis 1: Input Size vs. Execution Time}
Table \ref{tab:input_size_time} and Figure \ref{fig:input_size_time} show the empirical execution times as text length increases from $500$ to $10,000$ characters.

\begin{table}[H]
\centering
\caption{Empirical Execution Times (in milliseconds) across variable text sizes.}
\label{tab:input_size_time}
\vspace{0.2cm}
\resizebox{0.95\textwidth}{!}{
\begin{tabular}{ccccc}
\toprule
\textbf{Text Size ($n$)} & \textbf{Pattern Length ($m$)} & \textbf{Naive Match (ms)} & \textbf{KMP Match (ms)} & \textbf{Rabin-Karp (ms)} \\
\midrule
500   & 5  & 0.0456 & 0.0288 & 0.0621 \\
1000  & 10 & 0.0914 & 0.0597 & 0.1237 \\
10000 & 20 & 0.9339 & 0.5860 & 1.2857 \\
\bottomrule
\end{tabular}
}
\end{table}

\begin{figure}[H]
\centering
\includegraphics[width=0.88\textwidth]{images/input_size_vs_execution_time.png}
\caption{Input Size ($n$) versus Execution Time (ms) for Naive, KMP, and Rabin-Karp.}
\label{fig:input_size_time}
\end{figure}

\textbf{Key Insights}:
\begin{enumerate}
    \item \textbf{Linear Scaling}: All three algorithms exhibit linear scaling with respect to input text size $n$.
    \item \textbf{KMP Superiority}: KMP consistently records the lowest execution time ($0.586$ ms for $n = 10,000$), approximately $37\%$ faster than the Naive method and $54\%$ faster than Rabin-Karp. This is because KMP requires zero backtracking and minimal instruction cycles per character.
    \item \textbf{Rabin-Karp Overhead}: Although asymptotically $\mathcal{O}(n + m)$ on average, Rabin-Karp experiences higher constant-factor overhead in Python due to modular arithmetic (multiplication, subtraction, and modulo operations) at every window position.
\end{enumerate}

\newpage

\subsection*{4.4 Analysis 2: Pattern Length vs. Execution Time}
To isolate the effect of pattern length $m$ on runtime, the text length is held strictly constant at $n = 10,000$ characters while varying pattern length $m \in \{5, 10, 20\}$.

\begin{table}[H]
\centering
\caption{Execution Times (ms) across variable pattern lengths for fixed text length $n = 10,000$.}
\label{tab:pattern_length_time}
\vspace{0.2cm}
\resizebox{0.9\textwidth}{!}{
\begin{tabular}{cccc}
\toprule
\textbf{Pattern Length ($m$)} & \textbf{Naive Match (ms)} & \textbf{KMP Match (ms)} & \textbf{Rabin-Karp (ms)} \\
\midrule
5  & 1.0046 & 0.6008 & 1.4073 \\
10 & 1.0445 & 0.5801 & 1.3073 \\
20 & 0.9624 & 0.6036 & 1.5088 \\
\bottomrule
\end{tabular}
}
\end{table}

\begin{figure}[H]
\centering
\includegraphics[width=0.88\textwidth]{images/pattern_length_vs_execution_time.png}
\caption{Pattern Length ($m$) versus Execution Time (ms) with constant text size $n = 10,000$.}
\label{fig:pattern_length_time}
\end{figure}

\textbf{Key Insights}:
\begin{enumerate}
    \item \textbf{Invariance to Pattern Length}: When text size $n$ is held constant, varying pattern length $m$ from $5$ to $20$ produces virtually negligible impact on execution times across all three algorithms.
    \item \textbf{Why Naive Stays Flat}: In random text over $\Sigma = \{a,b,c,d\}$, mismatches occur within the first 1--2 characters regardless of whether the pattern has length 5 or 20.
    \item \textbf{KMP Stability}: The $\Theta(m)$ LPS preprocessing time ($m \le 20$) is dwarfed by the $\Theta(n)$ search time, maintaining flat line performance at $\approx 0.60$ ms.
    \item \textbf{Rabin-Karp Verification}: With $Q=101$, the number of collisions is very small, so the verification time remains bounded and independent of $m$.
\end{enumerate}

\newpage

\subsection*{4.5 Analysis 3: Input Size vs. Number of Comparisons}
Table \ref{tab:comparisons} and Figure \ref{fig:comparisons} detail the actual character comparisons executed by each algorithm as text size increases.

\begin{table}[H]
\centering
\caption{Number of Character Comparisons across different input text sizes.}
\label{tab:comparisons}
\vspace{0.2cm}
\resizebox{0.95\textwidth}{!}{
\begin{tabular}{ccccc}
\toprule
\textbf{Text Size ($n$)} & \textbf{Pattern ($m$)} & \textbf{Naive Comparisons} & \textbf{KMP Comparisons} & \textbf{Rabin-Karp Char Comps} \\
\midrule
500   & 5  & 658    & 627    & \textbf{21} \\
1000  & 10 & 1342   & 1257   & \textbf{20} \\
10000 & 20 & 13374  & 12561  & \textbf{187} \\
\bottomrule
\end{tabular}
}
\end{table}

\begin{figure}[H]
\centering
\includegraphics[width=0.88\textwidth]{images/input_size_vs_comparisons.png}
\caption{Input Size ($n$) versus Number of Character Comparisons.}
\label{fig:comparisons}
\end{figure}

\textbf{Key Insights}:
\begin{enumerate}
    \item \textbf{Rabin-Karp Comparison Efficiency}: Figure \ref{fig:comparisons} demonstrates the core algorithmic advantage of Rabin-Karp: it performs \textbf{only 187 character comparisons} for $n = 10,000$, compared to \textbf{13,374} for Naive and \textbf{12,561} for KMP. That is an astounding \textbf{98.6\% reduction} in character comparisons!
    \item \textbf{Why Rabin-Karp Performs So Few Character Comparisons}: Rabin-Karp filters almost all window positions via rolling integer arithmetic. Character equality checks occur exclusively when hash values coincide.
    \item \textbf{Naive versus KMP Comparisons}: KMP requires fewer comparisons than the Naive algorithm (12,561 vs 13,374 for $n=10,000$), saving comparisons because the LPS table skips re-evaluating known prefix alignments upon mismatch.
\end{enumerate}

\newpage

\subsection*{4.6 Comprehensive Algorithmic Comparison Matrix}
Table \ref{tab:comparison_matrix} summarizes the theoretical characteristics, empirical performance, and operational trade-offs of the three string matching algorithms.

\begin{table}[H]
\centering
\caption{Comprehensive Comparison of String Matching Algorithms.}
\label{tab:comparison_matrix}
\vspace{0.2cm}
\resizebox{\textwidth}{!}{
\begin{tabular}{lcccc}
\toprule
\textbf{Evaluation Metric} & \textbf{Naive Match} & \textbf{KMP Algorithm} & \textbf{Rabin-Karp Algorithm} \\
\midrule
\textbf{Best-Case Time Complexity}     & $\mathcal{O}(n)$ & $\Theta(n)$ & $\mathcal{O}(n)$ \\
\textbf{Average-Case Time Complexity}  & $\mathcal{O}(n)$ & $\Theta(n)$ & $\mathcal{O}(n)$ \\
\textbf{Worst-Case Time Complexity}    & $\mathcal{O}(n \cdot m)$ & $\Theta(n + m)$ & $\mathcal{O}(n \cdot m)$ \\
\textbf{Preprocessing Time}            & None ($\mathcal{O}(1)$) & $\Theta(m)$ (LPS Array) & $\Theta(m)$ (Initial Hashes) \\
\textbf{Auxiliary Space}               & $\mathcal{O}(1)$ & $\Theta(m)$ & $\mathcal{O}(1)$ \\
\textbf{Backtracking in Text}          & Yes & No (Strictly forward) & No (Sliding window) \\
\textbf{Character Comparisons ($n=10000$)} & High ($13,374$) & Moderate ($12,561$) & Extremely Low ($187$) \\
\textbf{Execution Speed (Empirical)}   & Fast ($0.93$ ms) & Fastest ($0.58$ ms) & Moderate ($1.28$ ms) \\
\textbf{Multiple Pattern Matching}     & Impractical & Requires Aho-Corasick & Natural extension (Bloom Filter) \\
\textbf{Best Use Case}                 & Small text / simple tasks & Large text / streaming data & Plagiarism / Multi-pattern \\
\bottomrule
\end{tabular}
}
\end{table}

\subsection*{4.7 Conclusion}
The laboratory experiments provide concrete empirical validation of theoretical string matching principles:
\begin{itemize}
    \item \textbf{KMP} offers the best overall execution performance and guarantees strict linear worst-case complexity $\Theta(n + m)$ without backtracking. It is the optimal algorithm for single-pattern search in continuous streams and large text corpora.
    \item \textbf{Rabin-Karp} dramatically reduces character comparisons by replacing string comparisons with $O(1)$ numerical hash operations, making it particularly valuable when character comparisons are computationally expensive or when extending to multi-pattern search.
    \item \textbf{Naive String Matching}, despite its worst-case $O(n \cdot m)$ vulnerability on periodic texts, performs adequately on average random text due to frequent early mismatches, while requiring zero preprocessing and $O(1)$ memory.
\end{itemize}

\end{document}
